\documentclass[11pt,a4paper]{amsart}
\usepackage{amsmath,amssymb,amsthm}
\usepackage{hyperref}

\newtheorem{theorem}{Theorem}[section]
\newtheorem{proposition}[theorem]{Proposition}
\newtheorem{lemma}[theorem]{Lemma}
\newtheorem{corollary}[theorem]{Corollary}
\newtheorem{definition}[theorem]{Definition}
\newtheorem{remark}[theorem]{Remark}

\title{P03 --- Haar-$L^2$ Firewall: $B_W$ Cannot Close on $L^2(\mathbb{R},du)$}
\author{Objekt-X-Programm}
\date{2026-08-08 (SYN, Entwurf)}

\begin{document}
\maketitle

\begin{abstract}
We consolidate NEU-253--258 and the NEU-257 Patch.
We show: (i) $L^2$-semiboundedness of $B_W \Leftrightarrow$ RH;
(ii) under RH, $B_W$ is not closable on Haar-$L^2(\mathbb{R},du)$;
(iii) Kato/KLMN on $L^2(\mathbb{R},du)$ fails ($\times[M]$).
We identify $\mathcal{H}_W \cong \ell^2(\Gamma,m_\gamma)$ under RH
(Suzuki 2025). The correct Weil Hilbert geometry is singular
relative to Lebesgue measure. All statements RH-conditional
as indicated.
\end{abstract}

\section{Background Space and Semiboundedness}

\begin{theorem}[$L^2$-Semiboundedness $\Leftrightarrow$ RH, NEU-257 $\checkmark$]
\[
  B_W \text{ is semibounded on } L^2(\mathbb{R},du) \iff \mathrm{RH}.
\]
\end{theorem}

\section{Non-Closability under RH}

\begin{theorem}[Non-closability, NEU-257 $\checkmark$]
Assuming RH,
\[
  B_W \text{ is not closable on } L^2(\mathbb{R},du).
\]
\end{theorem}

\begin{proof}
The spectral measure $\mu_W = \sum_\gamma m_\gamma\delta_\gamma$ is
singular with respect to Lebesgue measure; the Bochner--Schwartz
construction requires a translation-invariant measure compatible
with $\mu_W$, which $du$ is not. See NEU-257 Patch for full argument.
\end{proof}

\section{Kato/KLMN No-Go}

\begin{theorem}[Kato/KLMN $\times[M]$, NEU-257 $\checkmark$]
The Kato--Lions--Milgram--Necas theorem does not apply to
$B_W$ on $H_0=L^2(\mathbb{R},du)$.
\end{theorem}

\section{Weil Hilbert Space under RH}

\begin{theorem}[Weil Hilbert space, Suzuki 2025, NEU-257 $\checkmark$]
Assuming RH,
\[
  \mathcal{H}_W \cong L^2(\tau) \cong \ell^2(\Gamma,m_\gamma).
\]
\end{theorem}

\section{Consequence: Correct Geometry is Singular}

\begin{remark}
The correct Weil Hilbert geometry is not an operator on
$L^2(\mathbb{R},du)$; it is a singular spectral geometry
supported on zero-ordinates. This motivates the finite
interval construction of P04.
\end{remark}

\section{Source Nodes}

NEU-253, 254, 255, 256, 257 (Patch), 258
(archive-nodes/07-weil-explizitformel/).

\end{document}
